% GENERATED FILE. Edit canonical lessons, then run npm run generate:books.
% canonical-source-hash: fnv1a64:0ad65e25ccc79e46
% canonical-lessons: AR-C336-ahda, AR-C336-istaqbala, AR-C336-baraka, AR-C336-mahr, AR-C336-ashiyya

\chapter{Gifts and blessings}
\label{ch:ar-gifts-and-blessings}
\hlchaptermodality{336}

\begin{chapteropening}
\textbf{By the end of this chapter:} \emph{I can say to gift, to welcome, to congratulate, a dowry and the eve.}

Complete the last lesson of chapter 336: I can say to gift, to welcome, to congratulate, a dowry and the eve.
\end{chapteropening}

\section[\texorpdfstring{\ar{أهدى} (ahdā)}{أهدى (ahdā)}]{\ar{أهدى} (ahdā) — to gift; to dedicate}

\label{lesson:AR-C336-ahda}



\begin{quote}
Before the new one: say the Arabic for a giant, then the Arabic for a dwarf.
\end{quote}

\subsection*{You'll want to know: \ar{أهدى}}
\textbf{\ar{أهدى}} — \emph{ahdā} — \textquotedblleft{}to gift; to dedicate\textquotedblright{}.

\begin{quote}
\ar{أهدت} \ar{البنت} \ar{وردة} \ar{إلى} \ar{أمها}.
\end{quote}

\begin{quote}
\emph{ahdat al-bint wardatan ilā ummihā.}
\end{quote}

\begin{quote}
The girl gave her mother a rose.
\end{quote}

Present: \ar{يهدي}, \emph{yuhdī}.

The taught \ar{هدية}, a gift, is its noun. Of a book: \ar{أهدى} \ar{الكتاب} \ar{إلى} \ar{أمه}, he dedicated it to his mother.

\begin{cousinweb}
\textbf{\ar{أهدى}} (\emph{ahdā}) is built on the native Arabic root \emph{h-d-y} (Lane), \textquotedblleft{}to guide\textquotedblright{}, the root of the taught \emph{hadiyya}, a gift. It is Form IV, \textquotedblleft{}to send a gift\textquotedblright{}.
\end{cousinweb}

\begin{grammarlens}[title={What you've built}]
One more word for telling what happened.
\end{grammarlens}

\subsection*{Your turn}
\begin{itemize}
\raggedright
  \item \emph{Say it:} \emph{ahdā}
  \item \emph{Say it:} \emph{ahdā}, once more
  \item \emph{Say it:} \emph{ahdat al-bint wardatan ilā ummihā.} — The girl gave her mother a rose.
  \item \emph{Recall:} say the Arabic for a giant, then the Arabic for a dwarf, then say \emph{ahdā} again
\end{itemize}

\begin{tcolorbox}[breakable,colback=teal!4,colframe=teal!35!black,title={Before you move on}]
What does \ar{أهدى} mean? (\textquotedblleft{}To gift; to dedicate\textquotedblright{}.) Say it once more.
\end{tcolorbox}
\section[\texorpdfstring{\ar{استقبل} (istaqbala)}{استقبل (istaqbala)}]{\ar{استقبل} (istaqbala) — to receive guests; to welcome}

\label{lesson:AR-C336-istaqbala}



\begin{quote}
Before the new one: say the Arabic for a dwarf, then the Arabic for to gift.
\end{quote}

\subsection*{You'll want to know: \ar{استقبل}}
\textbf{\ar{استقبل}} — \emph{istaqbala} — \textquotedblleft{}to receive guests; to welcome\textquotedblright{}.

\begin{quote}
\ar{استقبل} \ar{الأهل} \ar{الضيوف} \ar{عند} \ar{الباب}.
\end{quote}

\begin{quote}
\emph{istaqbala al-ahl aḍ-ḍuyūf ʿinda al-bāb.}
\end{quote}

\begin{quote}
The family welcomed the guests at the door.
\end{quote}

Present: \ar{يستقبل}, \emph{yastaqbil}.

A hotel reception desk is \ar{الاستقبال}.

\begin{cousinweb}
\textbf{\ar{استقبل}} (\emph{istaqbala}) is built on the native Arabic root \emph{q-b-l} (Lane), \textquotedblleft{}to be in front of, to face\textquotedblright{}. It is Form X, \textquotedblleft{}to go to meet\textquotedblright{}; \emph{mustaqbal}, the future, is what comes to meet us.
\end{cousinweb}

\begin{grammarlens}[title={What you've built}]
One more word for telling what happened.
\end{grammarlens}

\subsection*{Your turn}
\begin{itemize}
\raggedright
  \item \emph{Say it:} \emph{istaqbala}
  \item \emph{Say it:} \emph{istaqbala}, once more
  \item \emph{Say it:} \emph{istaqbala al-ahl aḍ-ḍuyūf ʿinda al-bāb.} — The family welcomed the guests at the door.
  \item \emph{Recall:} say the Arabic for a dwarf, then the Arabic for to gift, then say \emph{istaqbala} again
\end{itemize}

\begin{tcolorbox}[breakable,colback=teal!4,colframe=teal!35!black,title={Before you move on}]
What does \ar{استقبل} mean? (\textquotedblleft{}To receive guests; to welcome\textquotedblright{}.) Say it once more.
\end{tcolorbox}
\section[\texorpdfstring{\ar{بارك} (bāraka)}{بارك (bāraka)}]{\ar{بارك} (bāraka) — to bless; to congratulate}

\label{lesson:AR-C336-baraka}



\begin{quote}
Before the new one: say the Arabic for to gift, then the Arabic for to receive guests.
\end{quote}

\subsection*{You'll want to know: \ar{بارك}}
\textbf{\ar{بارك}} — \emph{bāraka} — \textquotedblleft{}to bless; to congratulate\textquotedblright{}.

\begin{quote}
\ar{بارك} \ar{الأصدقاء} \ar{للعروسين}.
\end{quote}

\begin{quote}
\emph{bāraka al-aṣdiqāʾ li-l-ʿarūsayn.}
\end{quote}

\begin{quote}
The friends congratulated the bride and groom.
\end{quote}

Present: \ar{يبارك}, \emph{yubārik}.

\ar{بارك} \ar{لـ} someone: to congratulate them. The wish \ar{مبروك} is the same root.

\begin{cousinweb}
\textbf{\ar{بارك}} (\emph{bāraka}) is built on the native Arabic root \emph{b-r-k} (Lane), said of a camel that kneels and stays. \emph{Baraka}, a blessing, is good that stays; Form III \emph{bāraka} is \textquotedblleft{}to wish blessing on\textquotedblright{}.
\end{cousinweb}

\begin{grammarlens}[title={What you've built}]
One more word for telling what happened.
\end{grammarlens}

\subsection*{Your turn}
\begin{itemize}
\raggedright
  \item \emph{Say it:} \emph{bāraka}
  \item \emph{Say it:} \emph{bāraka}, once more
  \item \emph{Say it:} \emph{bāraka al-aṣdiqāʾ li-l-ʿarūsayn.} — The friends congratulated the bride and groom.
  \item \emph{Recall:} say the Arabic for to gift, then the Arabic for to receive guests, then say \emph{bāraka} again
\end{itemize}

\begin{tcolorbox}[breakable,colback=teal!4,colframe=teal!35!black,title={Before you move on}]
What does \ar{بارك} mean? (\textquotedblleft{}To bless; to congratulate\textquotedblright{}.) Say it once more.
\end{tcolorbox}
\section[\texorpdfstring{\ar{مهر} (mahr)}{مهر (mahr)}]{\ar{مهر} (mahr) — a dowry, a bride gift}

\label{lesson:AR-C336-mahr}



\begin{quote}
Before the new one: say the Arabic for to receive guests, then the Arabic for to bless.
\end{quote}

\subsection*{You'll want to know: \ar{مهر}}
\textbf{\ar{مهر}} — \emph{mahr} — \textquotedblleft{}a dowry, a bride gift\textquotedblright{}.

\begin{quote}
\ar{دفع} \ar{العريس} \ar{مهرا} \ar{للعروس}.
\end{quote}

\begin{quote}
\emph{dafaʿa al-ʿarīs mahran li-l-ʿarūs.}
\end{quote}

\begin{quote}
The groom paid a bride gift to the bride.
\end{quote}

The groom gives it to the bride; it is hers.

\begin{cousinweb}
\textbf{\ar{مهر}} (\emph{mahr}) is built on the native Arabic root \emph{m-h-r} (Lane), where \emph{mahr} is the gift a man gives his bride; \emph{muhr}, a foal, is written the same.
\end{cousinweb}

\begin{grammarlens}[title={What you've built}]
One more word for telling what happened.
\end{grammarlens}

\subsection*{Your turn}
\begin{itemize}
\raggedright
  \item \emph{Say it:} \emph{mahr}
  \item \emph{Say it:} \emph{mahr}, once more
  \item \emph{Say it:} \emph{dafaʿa al-ʿarīs mahran li-l-ʿarūs.} — The groom paid a bride gift to the bride.
  \item \emph{Recall:} say the Arabic for to receive guests, then the Arabic for to bless, then say \emph{mahr} again
\end{itemize}

\begin{tcolorbox}[breakable,colback=teal!4,colframe=teal!35!black,title={Before you move on}]
What does \ar{مهر} mean? (\textquotedblleft{}A dowry, a bride gift\textquotedblright{}.) Say it once more.
\end{tcolorbox}
\section[\texorpdfstring{\ar{عشية} (ʿashiyya)}{عشية (ʿashiyya)}]{\ar{عشية} (ʿashiyya) — the eve; the evening before}

\label{lesson:AR-C336-ashiyya}



\begin{quote}
Before the new one: say the Arabic for to bless, then the Arabic for a dowry.
\end{quote}

\subsection*{You'll want to know: \ar{عشية}}
\textbf{\ar{عشية}} — \emph{ʿashiyya} — \textquotedblleft{}the eve; the evening before\textquotedblright{}.

\begin{quote}
\ar{سافرنا} \ar{عشية} \ar{العيد}.
\end{quote}

\begin{quote}
\emph{sāfarnā ʿashiyyat al-ʿīd.}
\end{quote}

\begin{quote}
We travelled on the eve of the holiday.
\end{quote}

\ar{عشية} \ar{الانتخابات}: on the eve of the election.

\begin{cousinweb}
\textbf{\ar{عشية}} (\emph{ʿashiyya}) is built on the native Arabic root \emph{ʿ-sh-w} (Lane), where \emph{ʿashiyy} is the evening; the taught \emph{ʿashāʾ}, dinner, is the evening meal.
\end{cousinweb}

\begin{grammarlens}[title={What you've built}]
One more word for telling what happened.
\end{grammarlens}

\subsection*{Your turn}
\begin{itemize}
\raggedright
  \item \emph{Say it:} \emph{ʿashiyya}
  \item \emph{Say it:} \emph{ʿashiyya}, once more
  \item \emph{Say it:} \emph{sāfarnā ʿashiyyat al-ʿīd.} — We travelled on the eve of the holiday.
  \item \emph{Recall:} say the Arabic for to bless, then the Arabic for a dowry, then say \emph{ʿashiyya} again
\end{itemize}

\begin{tcolorbox}[breakable,colback=teal!4,colframe=teal!35!black,title={Before you move on}]
What does \ar{عشية} mean? (\textquotedblleft{}The eve; the evening before\textquotedblright{}.) Say it once more.
\end{tcolorbox}
